\documentclass[10pt,a4paper,twocolumn]{article}
\usepackage[utf8]{inputenc}
\usepackage[T5]{fontenc} 
\usepackage{amsmath}
\usepackage{graphicx}
\usepackage{hyperref}
\usepackage{geometry}
\usepackage{booktabs}
\geometry{a4paper, left=0.75in, right=0.75in, top=1in, bottom=1in, columnsep=0.25in}

\title{ECG Heartbeat Classification: A Deep Transferable Representation\\ \large Project Report}
\author{Nguyễn Thị Như Quỳnh}
\date{}

\begin{document}
\maketitle

\section{Dataset Description}
The dataset utilized in this project is extracted from the MIT-BIH Arrhythmia Database, carefully preprocessed and formatted for machine learning tasks. It consists of 1D ECG heartbeat signals, categorized into five distinct classes based on the AAMI EC57 standard:
\begin{itemize}
    \item \textbf{Class 0 (N):} Normal beat
    \item \textbf{Class 1 (S):} Supraventricular ectopic beat
    \item \textbf{Class 2 (V):} Ventricular ectopic beat
    \item \textbf{Class 3 (F):} Fusion beat
    \item \textbf{Class 4 (Q):} Unknown beat
\end{itemize}

Figure \ref{fig:ecg_samples} illustrates a sample waveform from each of these categories. To ensure consistent input dimensions for the convolutional neural network, all signals were zero-padded to a uniform length of 187 time steps (visible as the flat line at the tail of each signal). 

\begin{figure}[h]
    \centering
    \includegraphics[width=\linewidth]{tải xuống (48).png}
    \caption{Sample 1D ECG waveforms for each of the five heartbeat classes.}
    \label{fig:ecg_samples}
\end{figure}

The most prominent challenge of this dataset is its severe class imbalance. The training set contains 87,554 samples, with Class 0 (Normal) heavily dominating the distribution (72,471 samples, roughly 82.7\%), while minority classes like Fusion beats constitute a very small fraction (641 samples).

\section{Model Implementation}
To capture the spatial and temporal characteristics of the ECG signals, a 1-Dimensional Convolutional Neural Network (1D CNN) was constructed. 

The network architecture consists of two main convolutional blocks. The first block applies 64 filters of size 5, followed by Batch Normalization to stabilize the activations, and a MaxPooling1D layer (pool size 2) for spatial dimensionality reduction. The second block applies 128 filters of size 3, again followed by Batch Normalization and MaxPooling1D.

The extracted feature maps are then flattened and passed through a Dense layer of 128 neurons with a ReLU activation function. To prevent overfitting, a Dropout layer with a rate of 0.5 is applied. The final output is generated by a Dense layer with 5 neurons using a \textit{softmax} activation function, providing the probability distribution across the five heartbeat classes. The model was compiled using the Adam optimizer and categorical cross-entropy loss function.

\section{Results and Hyperparameter Experiments}
During the initial training phase, a batch size of 64 was used over 10 epochs. To counter the severe dataset imbalance, balanced class weights were applied directly to the loss function. 

\subsection{Performance Analysis}
The baseline experimental run yielded an overall test accuracy of 82.76\%. However, a detailed examination of the classification metrics revealed significant challenges regarding model generalization:

\begin{table}[h]
\centering
\begin{tabular}{lccc}
\toprule
\textbf{Class} & \textbf{Precision} & \textbf{Recall} & \textbf{F1-Score} \\
\midrule
Normal & 0.83 & 1.00 & 0.91 \\
Supravent. & 0.00 & 0.00 & 0.00 \\
Ventricular & 0.00 & 0.00 & 0.00 \\
Fusion & 0.00 & 0.00 & 0.00 \\
Unknown & 0.00 & 0.00 & 0.00 \\
\bottomrule
\end{tabular}
\caption{Classification Report on Test Set}
\label{tab:results}
\end{table}

The training log indicates that the model achieved a 1.0000 (100\%) accuracy on the training set by Epoch 3, while the validation loss exploded exponentially (from 85.87 to 457.52). The F1-scores of 0.00 for all minority classes confirm that the model experienced severe overfitting and gradient instability, ultimately collapsing into a zero-rule predictor that blindly guesses the majority class ("Normal").

\subsection{Hyperparameter Optimization Strategy}
Based on these empirical results, applying static class weights combined with Batch Normalization in this specific CNN architecture causes internal covariate shifts that hinder the learning of minority features. To improve performance in future iterations, the following hyperparameter adjustments should be experimented with:
\begin{itemize}
    \item \textbf{Learning Rate:} Reducing the Adam optimizer's learning rate from the default 0.001 to 0.0001 to prevent the loss from diverging and exploding during weight updates.
    \item \textbf{Dropout Tuning:} Increasing the dropout rate or adding L2 regularization to further constrain the network and prevent it from memorizing the training data.
    \item \textbf{Alternative Sampling:} Removing the `class\_weight` parameter and instead utilizing SMOTE (Synthetic Minority Over-sampling Technique) to physically balance the dataset before feeding it into the neural network.
\end{itemize}

\end{document}